%% kapa: exam=1 round=31 year=2020 period=2 form=A subject=accounting subject_ko=회계학 seq=3 number=43
%% source: Q-Net 감정평가사 기출문제 (cst00309) artlSeq=5207219 file=extracted/1008/5207219/A형/A형/2020년도 제31회 감정평가사 1차 시험_2교시_A형.pdf
%% conversion: kordoc 4.24.51 -> tools/gen1.py (auto); review=proofread-sonnet; flags=none
\begin{question}{43}
(주)감평은 20x1년 1월 1일 미국에 있는 건물(취득원가 \$5,000, 내용연수 5년, 잔존가치 \$0, 정액법 상각)을 취득하였다. (주)감평은 건물에 대하여 재평가모형을 적용하고 있으며, 20x1년 12월 31일 현재 동 건물의 공정가치는 \$6,000로 장부금액과의 차이는 중요하다. (주)감평의 기능통화는 원화이며, 20x1년 1월 1일과 20x1년 12월 31일의 환율은 각각 ￦1,800/\$과 ￦1,500/\$이고, 20x1년의 평균환율은 ￦1,650/\$이다. (주)감평이 20x1년 말 재무상태표에 인식해야 할 건물에 대한 재평가잉여금은?
\begin{choices}
  \item ￦1,500,000
  \item ￦1,650,000
  \item ￦1,800,000
  \item ￦3,000,000
  \item ￦3,300,000
\end{choices}
\end{question}
\answer{3}
